\chapter{Introduction}
Modern warfare has transitioned into an increasingly data-driven environment where split-second decisions are critical for survival and mission success. The sheer volume of visual information from surveillance feeds, drones, and tactical environments often overwhelms personnel, leading to cognitive fatigue and a higher margin for error. Manual threat assessment is no longer sufficient to ensure the safety and tactical superiority of ground units in high-stress scenarios. Consequently, there is an urgent need for intelligent, autonomous systems that can process vast amounts of real-time visual data to provide actionable insights and enhance situational awareness. \\
VeerDrishti is an AI-enabled combat vision system designed to bridge this gap by providing soldiers with a persistent, intelligent observer through a camera-integrated solution. By leveraging state-of-the-art computer vision models and edge computing, the system autonomously identifies enemy combatants vehicles, and weapons. The primary objective is to offload the mental burden of constant monitoring from the soldier, allowing them to focus on strategic execution. Developed as a lightweight solution suitable for deployment on wearable hardware, the system aims to transform how visual data is utilized on the front lines of national security. \\
The project incorporates sophisticated methodologies to ensure reliability under diverse and challenging battlefield conditions. It utilizes multimodal image fusion combining visible and thermal spectrums—to maintain high detection accuracy in low-light or obscured environments. Furthermore, by implementing edge-optimized architectures and advanced feature fusion techniques, the system achieves high-speed inference without compromising precision. This integration of versatile object detection and robust information processing ensures that the system remains a reliable asset, offering a significant technological advantage in modern combat and defense operations.

\section{Problem Statement}
The problem statement is to address the challenges of managing vast real-time visual data in modern battlefield environments, where manual monitoring is slow, error-prone, and inefficient; in order to enhance situational awareness and tactical decision-making, an AI-powered combat vision system is proposed to automatically detect, classify, and assess threats in real time, improving response speed and operational efficiency.

\section{Motivation}
The motivation for VeerDrishti is to minimize casualties and save lives by providing an intelligent, persistent observer in combat environments. Modern warfare involves overwhelming visual data that leads to cognitive fatigue, increasing the risk of fatal human error. The system serves as a technological safeguard, enhancing personnel safety and protecting national security forces during high-stress operations.

\section{Objectives}
The project's primary objectives are to develop a lightweight, real-time combat vision system that can be deployed on wearable hardware, capable of autonomously detecting and classifying threats in dynamic battlefield environments. Specific goals include: \\
i. Implement a lightweight real-time object detection model. \\
ii. Enable on-screen visual highlighting of identified threats. \\
iii. Optimize system performance for low latency and smooth operation. \\
iv. Evaluate accuracy and response time under simulated field conditions.

\section{Feasibility Study}
The feasibility of the proposed combat vision system is assessed across technical, operational, and economic dimensions.

\subsection{Technical Feasibility}
The proposed system is technically feasible using existing smartphone hardware and computer vision frameworks. Modern phones provide high-resolution cameras, sufficient processing power, and GPU support for real-time object detection. Lightweight AI models can be optimized for edge deployment, enabling on-device threat detection without requiring complex external hardware.

\subsection{Operational Feasibility}
The system can be integrated into existing helmet setups with minimal modification, using a mounted smartphone as both camera and display interface. It reduces cognitive load by automating threat detection and visual highlighting. Basic training would enable personnel to interpret AR overlays effectively in dynamic environments.

\subsection{Economic Feasibility}
Using commercially available smartphones significantly lowers development and deployment costs compared to specialized military hardware. Open-source AI frameworks further reduce software expenses. The prototype approach minimizes initial investment, making the system cost-effective for research, testing, and potential scalable improvements in future versions.

\section{Work Distribution}
\begin{table}[h]
    \centering
    \caption{Work Distribution Among Team Members}
    \begin{tabular}{lll}
        \toprule
        \textbf{Name}  & \textbf{Tasks}                       \\
        \midrule
        Gaurav Chauhan & Literature Survey, Report Writing    \\
        Gaurav Singh   & Literature Survey, Slide Preparation \\
        \bottomrule
    \end{tabular}
\end{table}